\documentclass[aspectratio=169,spanish,xcolor=table]{beamer}
\usepackage{iftex}
\ifPDFTeX
  \usepackage[utf8]{inputenc}
  \usepackage[T1]{fontenc}
\fi
\usepackage[spanish,es-tabla,es-noshorthands]{babel}
\usepackage{helvet}
\renewcommand{\familydefault}{\sfdefault}
\usepackage{tikz,pgfplots,array,booktabs,ragged2e}
\pgfplotsset{compat=1.18}
\usetikzlibrary{arrows.meta,positioning,calc,shapes.geometric,fit}

\definecolor{iiiepeNavy}{HTML}{163B5C}
\definecolor{iiiepeBlue}{HTML}{1F6F9F}
\definecolor{iiiepeAqua}{HTML}{20B7A8}
\definecolor{iiiepeOrange}{HTML}{F28C28}
\definecolor{iiiepePaper}{HTML}{F6F8F7}
\definecolor{iiiepeInk}{HTML}{1E2A32}
\definecolor{iiiepePale}{HTML}{EEF5F8}
\definecolor{iiiepeMint}{HTML}{E7F6F4}
\definecolor{iiiepeWarm}{HTML}{FFF3E6}

\usetheme{default}
\usecolortheme{default}
\setbeamertemplate{navigation symbols}{}
\setbeamercolor{normal text}{fg=iiiepeInk,bg=white}
\setbeamercolor{structure}{fg=iiiepeNavy}
\setbeamercolor{frametitle}{fg=white,bg=iiiepeNavy}
\setbeamercolor{block title}{fg=white,bg=iiiepeBlue}
\setbeamercolor{block body}{fg=iiiepeInk,bg=iiiepePale}
\setbeamercolor{item}{fg=iiiepeOrange}

\setbeamertemplate{frametitle}{%
  \nointerlineskip%
  \begin{beamercolorbox}[wd=\paperwidth,ht=0.95cm,dp=0.22cm,leftskip=0.45cm,rightskip=0.45cm]{frametitle}%
    \usebeamerfont{frametitle}\insertframetitle%
    \hfill\includegraphics[height=0.58cm]{img/iiiepe/logo-iiiepe-blanco.png}%
  \end{beamercolorbox}%
}

\setbeamertemplate{footline}{%
  \leavevmode%
  \hbox{%
    \begin{beamercolorbox}[wd=.78\paperwidth,ht=0.5cm,dp=0.18cm,leftskip=0.45cm]{frametitle}%
      \scriptsize\color{white}M. de la Cruz Muñoz · MAEM · Matrícula 00874 · SD26 · Actividad 2 · Sesión 2%
    \end{beamercolorbox}%
    \begin{beamercolorbox}[wd=.22\paperwidth,ht=0.5cm,dp=0.18cm,rightskip=0.45cm]{frametitle}%
      \hfill\scriptsize\color{white}\insertframenumber{} / \inserttotalframenumber%
    \end{beamercolorbox}%
  }%
  \vskip0pt%
}

\newcommand{\Pill}[2]{%
  \tikz[baseline=(x.base)]{\node[rounded corners=3pt,fill=#1!14,draw=#1!75!black,thick,inner xsep=6pt,inner ysep=4pt,align=center] (x) {\scriptsize\bfseries #2};}%
}

\title{¿Cómo Soy Como Maestro? Silueta Docente}
\subtitle{Actividad 2 · Sesión 2 · Módulo 50158}
\author{Mtro. Martín Jonathan de la Cruz Muñoz · Matrícula 00874}
\institute{IIIEPE · Maestría en Aprendizaje y Enseñanza de las Matemáticas · Grupo A}
\date{19 de noviembre de 2026}

\begin{document}

\begin{frame}[plain]
  \begin{tikzpicture}[remember picture,overlay]
    \fill[iiiepeNavy] (current page.north west) rectangle ([xshift=10.3cm]current page.south west);
    \fill[iiiepeBlue] ([xshift=10.3cm]current page.north west) rectangle (current page.south east);
    \fill[iiiepeOrange] ([xshift=0.8cm,yshift=-6.45cm]current page.north west) rectangle ([xshift=9.3cm,yshift=-6.53cm]current page.north west);
    \node[anchor=north west,text width=8.7cm,align=left] at ([xshift=0.8cm,yshift=-1.15cm]current page.north west) {%
      {\color{iiiepeAqua!45!white}\small\bfseries MAESTRÍA EN APRENDIZAJE Y ENSEÑANZA DE LAS MATEMÁTICAS · SD26\par}%
      \vspace{0.35cm}
      {\color{white}\Large\bfseries \inserttitle\par}%
      \vspace{0.25cm}
      {\color{white!90}\normalsize \insertsubtitle\par}%
    };
    \node[anchor=south west,text width=8.7cm,align=left] at ([xshift=0.8cm,yshift=0.75cm]current page.south west) {%
      {\color{white}\normalsize\bfseries Mtro. Martín Jonathan de la Cruz Muñoz\par}%
      {\color{white!85}\scriptsize Matrícula 00874 · martin.delacruz@campus.iiiepe.edu.mx\par}%
      {\color{iiiepeAqua!35!white}\scriptsize Monterrey, Nuevo León · 19 de noviembre de 2026\par}%
    };
    \node[anchor=center,text width=5.0cm,align=center] at ([xshift=13.15cm,yshift=-4.5cm]current page.north west) {%
      \includegraphics[width=3.6cm]{img/iiiepe/logo-iiiepe-blanco.png}\par\vspace{0.45cm}%
      {\color{white}\large\bfseries IIIEPE\par}%
      \vspace{0.1cm}%
      {\color{white!85}\scriptsize Fundamentos para la enseñanza\\y el aprendizaje I\par}%
    };
  \end{tikzpicture}
\end{frame}
\begin{frame}{Propósito de la Sesión 2: ¿Cómo Soy Como Maestro?}
  \begin{block}{Tesis de identidad docente}
    La docencia matemática madura articula la precisión técnica heredada de la ingeniería con la sensibilidad humanista, el diálogo respetuoso y la mediación equitativa \textbf{(Medina Rivilla y Salvador Mata, 2009; UNESCO, 2015)}.
  \end{block}
  \vspace{0.3cm}
  \begin{columns}[T]
    \column{0.48\textwidth}
      \textbf{\color{iiiepeNavy}Contexto de práctica:}
      \begin{itemize}\small
        \item Preparatoria técnica (CONALEP y CECyTE: mantenimiento y mecatrónica).
        \item Docencia en ingeniería (Mecánica, Sistemas e Industrial).
      \end{itemize}
    \column{0.48\textwidth}
      \textbf{\color{iiiepeBlue}Ejes del ejercicio reflexivo:}
      \begin{itemize}\small
        \item Cinco preguntas de autoindagación docente.
        \item Silueta docente: centro, lado izquierdo y lado derecho.
        \item Compromisos situados de mejora continua.
      \end{itemize}
  \end{columns}
\end{frame}

\begin{frame}{Cinco Preguntas de Reflexión Inicial (1--3)}
  \begin{itemize}\small
    \item \textbf{\color{iiiepeNavy}1. ¿Cómo describes tu forma de enseñar?} En transición desde el orden procedimental ingenieril hacia una mediación dialógica centrada en comprender por qué funciona un modelo.
    \vspace{0.2cm}
    \item \textbf{\color{iiiepeBlue}2. ¿Qué estrategias utilizas con mayor frecuencia?} Situaciones problema de taller, verificación dimensional, contraste tabla--gráfica--álgebra y discusión en equipos \textbf{(Pimienta Prieto, 2012)}.
    \vspace{0.2cm}
    \item \textbf{\color{iiiepeAqua!75!black}3. Principales fortalezas:} Dominio aplicado en contextos industriales, diseño visual estructurado y trato respetuoso hacia cada estudiante.
  \end{itemize}
\end{frame}

\begin{frame}{Cinco Preguntas de Reflexión Inicial (4--5)}
  \begin{columns}[T]
    \column{0.48\textwidth}
      \begin{block}{4. Aspectos por mejorar}
        \begin{itemize}\scriptsize
          \item Esperar más tiempo antes de resolver la dificultad en el pizarrón para favorecer la autonomía estudiantil.
          \item Incorporar evaluación formativa continua con criterios claros desde el inicio.
          \item Escalonar andamiajes ante trayectorias matemáticas heterogéneas \textbf{(Sawyer, 2008)}.
        \end{itemize}
      \end{block}
    \column{0.48\textwidth}
      \begin{block}{5. Toma de decisiones en clase}
        \begin{itemize}\scriptsize
          \item Pasar de intuición apresurada («repetir otro ejercicio») a diagnóstico conceptual del error.
          \item Usar evidencias breves de desempeño antes de cerrar un tema \textbf{(Collin y Quigley, 2021)}.
          \item Ajustar explicaciones según preguntas reales del grupo.
        \end{itemize}
      \end{block}
  \end{columns}
\end{frame}

\begin{frame}{Silueta Docente: Esquema Integrador}
  \centering
  \begin{tikzpicture}[font=\tiny, scale=0.80, transform shape]
    \draw[fill=iiiepePale,draw=iiiepeNavy,very thick] (0,2.55) circle (0.55cm);
    \node[align=center,font=\tiny\bfseries,iiiepeNavy] at (0,2.55) {RIGOR\\Y ÉTICA};
    \draw[fill=iiiepeBlue!10,draw=iiiepeNavy,very thick,rounded corners=8pt]
      (-0.25,2.0) -- (0.25,2.0) -- (0.95,1.72) -- (0.85,0.45) -- (0.65,0.45) -- (0.60,1.35)
      -- (0.55,-1.25) -- (0.12,-1.25) -- (0.04,-0.08) -- (-0.04,-0.08) -- (-0.12,-1.25) -- (-0.55,-1.25)
      -- (-0.60,1.35) -- (-0.65,0.45) -- (-0.85,0.45) -- (-0.95,1.72) -- cycle;
    \node[draw=iiiepeNavy,fill=white,rounded corners=2pt,text width=1.65cm,align=center,inner sep=3pt] at (0,1.25)
      {\textbf{CENTRO}\\Rigor · Diálogo · Dignidad};

    \node[draw=iiiepeNavy,fill=iiiepePale,rounded corners=3pt,text width=5.0cm,align=left,inner sep=5pt,anchor=north east] (l1) at (-1.35,2.9)
      {\textbf{\color{iiiepeNavy}IZQUIERDO · Fortalezas situadas:}\\
       Modelación en taller (calibración, costos) y análisis dimensional con respeto a las preguntas del grupo.};
    \node[draw=iiiepeOrange!90!black,fill=iiiepeWarm,rounded corners=3pt,text width=5.0cm,align=left,inner sep=5pt,anchor=north east] (l2) at (-1.35,0.6)
      {\textbf{\color{iiiepeOrange!85!black}IZQUIERDO · Áreas de oportunidad:}\\
       Resolver demasiado pronto por presión de tiempo; fortalecer retroalimentación formativa continua.};

    \node[draw=iiiepeBlue,fill=iiiepePale,rounded corners=3pt,text width=5.0cm,align=left,inner sep=5pt,anchor=north west] (r1) at (1.35,2.9)
      {\textbf{\color{iiiepeBlue}DERECHO · Necesidades de formación:}\\
       Andamiaje para aprendizaje profundo, evaluación formativa e investigación-acción docente.};
    \node[draw=iiiepeAqua!80!black,fill=iiiepeMint,rounded corners=3pt,text width=5.0cm,align=left,inner sep=5pt,anchor=north west] (r2) at (1.35,0.6)
      {\textbf{\color{iiiepeAqua!75!black}DERECHO · Bibliografía por explorar:}\\
       Sawyer (2008), Meirieu (2020), Restrepo Gómez (2004), Bronfenbrenner y Morris (2006).};

    \draw[->,thick,iiiepeNavy] (l1.east) -- (-0.75,1.65);
    \draw[->,thick,iiiepeOrange!90!black] (l2.east) -- (-0.58,-0.25);
    \draw[->,thick,iiiepeBlue] (0.75,1.65) -- (r1.west);
    \draw[->,thick,iiiepeAqua!80!black] (0.58,-0.25) -- (r2.west);
  \end{tikzpicture}
\end{frame}

\begin{frame}{Centro de la Silueta: Quién Soy Como Docente}
  \begin{block}{Síntesis de identidad profesional}
    «Docente de perfil técnico y de ingeniería en tránsito hacia una práctica humanista que concibe el rigor matemático como herramienta de comprensión compartida y equidad.»
  \end{block}
  \vspace{0.35cm}
  \centering
  \Pill{iiiepeNavy}{Rigor conceptual} \hspace{0.3cm}
  \Pill{iiiepeBlue}{Diálogo horizontal} \hspace{0.3cm}
  \Pill{iiiepeAqua!75!black}{Dignidad e inclusión}
  \vspace{0.35cm}
  \begin{itemize}\small
    \item Respeto profundo por el proceso intelectual de jóvenes de bachillerato técnico y adultos en ingeniería.
    \item Convicción de que el error es información didáctica valiosa y no motivo de exclusión.
  \end{itemize}
\end{frame}

\begin{frame}{Lado Izquierdo: Fortalezas y Áreas de Oportunidad Situadas}
  \scriptsize\renewcommand{\arraystretch}{1.35}
  \begin{tabular}{p{3.0cm}p{5.0cm}p{5.0cm}}
    \toprule
    \rowcolor{iiiepePale} \textbf{Eje} & \textbf{Situación real de aula} & \textbf{Lectura pedagógica} \\
    \midrule
    Fortaleza: Vinculación técnica & Lectura de gráficas y modelos lineales con sensores reales o tarifas de mantenimiento. & Otorga significado tangible a pendientes e intersecciones. \\
    Fortaleza: Escucha respetuosa & Análisis colectivo de errores dimensionales en pizarrón sin ridiculizar. & Disminuye ansiedad y habilita participación. \\
    Oportunidad: Cesión de control & Intervenir anticipadamente en el despeje por premura del horario escolar. & Limita la construcción autónoma de argumentos \textbf{(Sawyer, 2008)}. \\
    Oportunidad: Seguimiento continuo & Concentrar evaluación en exámenes de cierre de unidad. & Demanda evidencias diagnósticas frecuentes \textbf{(Boström y Palm, 2023)}. \\
    \bottomrule
  \end{tabular}
\end{frame}

\begin{frame}{Lado Derecho: Necesidades de Formación y Bibliografía}
  \begin{columns}[T]
    \column{0.48\textwidth}
      \begin{block}{Necesidades prioritarias}
        \begin{itemize}\scriptsize
          \item \textbf{Andamiaje sin suplantación:} Aprender a formular preguntas detonadoras que mantengan el reto cognitivo.
          \item \textbf{Evaluación para aprender:} Incorporar devoluciones formativas antes de calificar \textbf{(Collin y Quigley, 2021)}.
          \item \textbf{Indagación situada:} Sistematizar diarios y producciones estudiantiles como evidencia docente.
        \end{itemize}
      \end{block}
    \column{0.48\textwidth}
      \begin{block}{Autores que orientan mi ruta}
        \begin{itemize}\scriptsize
          \item \textbf{Sawyer (2008):} Ciencias del aprendizaje profundo.
          \item \textbf{Meirieu (2020):} Sentido colectivo y democrático del aula.
          \item \textbf{Restrepo Gómez (2004):} Investigación-acción pedagógica.
          \item \textbf{Bronfenbrenner y Morris (2006):} Contextos y procesos proximales.
        \end{itemize}
      \end{block}
  \end{columns}
\end{frame}

\begin{frame}{Compromisos Concretos en la Práctica Matemática}
  \begin{itemize}\normalsize
    \item \textbf{Regla de espera dialógica:} Reservar tiempo real para que equipos formulen conjeturas antes de validar procedimientos en el pizarrón.
    \item \textbf{Criterios transparentes:} Socializar metas de argumentación matemática desde el inicio de cada secuencia.
    \item \textbf{Andamiaje escalonado:} Ofrecer representaciones concretas y preguntas guía a quienes llegan con rezago operativo sin rebajar la meta conceptual.
    \item \textbf{Registro reflexivo:} Documentar obstáculos frecuentes de clase para investigar mi propia enseñanza.
  \end{itemize}
\end{frame}

\begin{frame}{Conclusiones de la Sesión 2}
  \begin{itemize}\normalsize
    \item La reflexión crítica revela que nuestras creencias implícitas gobiernan las decisiones rápidas dentro del aula.
    \item Mi fortaleza en el dominio técnico e ingenieril se potencia cuando se integra con escucha activa, retroalimentación descriptiva y equidad.
    \item La silueta docente funciona como línea base evaluable para contrastar avances al cierre de la cohorte SD26.
  \end{itemize}
\end{frame}

\begin{frame}{Referencias Principales (APA 7)}
  \scriptsize\setlength{\parskip}{5pt}
  \begin{itemize}
    \item Boström, E., y Palm, T. (2023). Formative assessment practice in mathematics. \textit{Frontiers in Education}, 8, 1101192.
    \item Bronfenbrenner, U., y Morris, P. A. (2006). The bioecological model of human development. En \textit{Handbook of child psychology} (pp. 793--828). Wiley.
    \item Collin, J., y Quigley, A. (2021). \textit{Teacher feedback to improve pupil learning}. Education Endowment Foundation.
    \item Medina Rivilla, A., y Salvador Mata, F. (2009). \textit{Didáctica general} (2.ª ed.). Pearson.
    \item Meirieu, P. (2020). ``La escuela después''... ¿Con la pedagogía de antes? \textit{MCEP}.
    \item Restrepo Gómez, B. (2004). La investigación-acción educativa y la construcción de saber pedagógico. \textit{Educación y Educadores}, (7), 45--55.
    \item Sawyer, R. K. (2008). Optimizing learning: Implications of learning sciences research. OECD.
  \end{itemize}
\end{frame}

\end{document}
